\begin{lemma}[Existence of the ordered union of two tree-sequences]
Let $A$ and $B$ be finite ordered tree-sequences in the same loopless
multigraph and with the same distinguished first edge identity $e_0$.
There is a tree-sequence $U$ such that
\[
                         \noderef{OrderedTreeUnion}(A,B,U).
\]
Thus $U$ is obtained from $A$ by appending precisely the edge--new-vertex
entries of $B$ whose new vertices lie outside $V(A)$, in their order in
$B$.
\end{lemma}
\begin{proof}
Write
\[
 A=(a_0,e_1,a_1,\ldots,e_p,a_p),\qquad
 B=(b_0,f_1,b_1,\ldots,f_q,b_q).
\]
By \noderef{TreeSequence}, $e_1=f_1=e_0$, and the attachment condition at
the first edge gives
\[
       \operatorname{ends}(e_0)=\{a_1,a_0\}=\{b_1,b_0\}.       \tag{1}
\]
Consequently $b_0,b_1\in V(A)$.  This observation also handles the
possibility that the two sequences orient the ends of $e_0$ oppositely.

Let
\[
 I=\{\,i\in\{1,\ldots,q\}:b_i\notin V(A)\,\},
\]
and list its members in their inherited order as
$i_1<\cdots<i_m$.  Starting with the complete displayed sequence $A$,
append, for $r=1,\ldots,m$, the exact entry $f_{i_r},b_{i_r}$ from $B$.
Call the resulting displayed sequence $U$.  We first verify that this is
a tree-sequence in the sense of \noderef{TreeSequence}.

The old vertices are pairwise distinct because they are the vertices of
$A$.  Each appended vertex lies outside $V(A)$ by the definition of $I$,
and two appended vertices are distinct because the vertex map of $B$ is
injective.  Hence all vertices displayed in $U$ are pairwise distinct.
The old edge identities are likewise pairwise distinct.  The appended
edge identities are pairwise distinct by edge injectivity in $B$.  No
appended edge can equal an old edge: the attachment clause for $A$ shows
that both ends of every old edge belong to $V(A)$, whereas the appended
edge $f_{i_r}$ has the end $b_{i_r}\notin V(A)$.  Thus the edge entries of
$U$ are also pairwise distinct.  Notice that this is an argument about
exact edge identities; it does not identify parallel edges having the
same unordered pair of ends.

It remains to check attachment for an appended entry $f_{i_r},b_{i_r}$.
The attachment clause for $B$ supplies an index $t<i_r$ such that
\[
                  \operatorname{ends}(f_{i_r})
                    =\{b_{i_r},b_t\}.                           \tag{2}
\]
If $b_t\in V(A)$, then \noderef{TreeSequenceVertexSet} supplies an old
position of $A$, hence of $U$, occupied by $b_t$, so (2) attaches the new
entry there.
If $b_t\notin V(A)$, then (1) implies $t\geq2$; hence $t$ is one of the
indices in $I$.  Since $t<i_r$, it occurs among
$i_1,\ldots,i_{r-1}$, so $b_t$ has already been appended to $U$ when the
$r$-th entry is added.  Equation (2) again gives the required attachment.
The first edge of $U$ is still the exact edge identity $e_0$, and its
length is positive because the initial segment $A$ has positive length.
This proves that $U$ is a tree-sequence.

We now check every field of \noderef{OrderedTreeUnion}.  By construction,
$A$ is an exact ordered initial segment of $U$: all of its vertex and edge
entries are unchanged.  The vertices added after $A$ are exactly the
$b_i$ outside $V(A)$, while every vertex of $B$ not outside $V(A)$ already
belongs to $V(A)$.  Using \noderef{TreeSequenceVertexSet}, this proves
\[
                         V(U)=V(A)\cup V(B).                     \tag{3}
\]
Every entry after $A$ is, by construction, the exact pair
$f_{i_r},b_{i_r}$ from $B$, which proves the new-entry clause.  Conversely,
if a new vertex $b_i$ of an entry of $B$ lies outside $V(A)$, then
$i\in I$, so that exact entry occurs in the appended part of $U$; this is
the coverage clause.  Finally, if $i<j$ are two such indices, their ranks
in the increasing list $i_1<\cdots<i_m$ have the same strict order.
Their positions in $U$ therefore have that strict order, proving the
right-order clause.  These facts together are precisely an instance of
\noderef{OrderedTreeUnion}$(A,B,U)$.
\end{proof}
